\chapter{Compter ce que coûte le calcul}\label{ch:sobriete}

Un assistant consomme de l'électricité, et les machines qui le font tourner ont été
fabriquées. On parle d'IA frugale, ou de sobriété numérique, pour l'idée simple qu'il faut
le savoir avant de s'en servir beaucoup, et s'en servir moins quand on peut. Ce projet a
fait le compte en novembre 2026, pour ses trois premiers mois (du 7 août au 7 novembre
2026). Le détail, avec les sources et les calculs, est dans l'article sur le coût carbone
du projet (\fichier{docs/carbon/article.tex}). On n'en reprend ici que la méthode et les
ordres de grandeur, puis les gestes qui en découlent.

\section{Où part l'énergie}

Un modèle de langage dépense de l'énergie de deux façons. Il lit le texte qu'on lui
soumet, d'un bloc, puis il écrit sa réponse, jeton après jeton. Pour ne pas tout relire, il
garde de chaque conversation un résumé en mémoire, le cache, qu'il récupère au tour
suivant. Une lecture depuis le cache coûte bien moins qu'une lecture neuve, mais elle n'est
pas gratuite.

Or un agent comme celui de ce livre renvoie au modèle, à chaque tour, toute la conversation
depuis le début : les fichiers lus, les résultats des commandes, les réponses déjà
données. Du 24 septembre au 7 novembre 2026, chaque appel au modèle principal a relu en
moyenne environ 515\,000 jetons depuis le cache, en a écrit environ 5\,200 dans le cache
et a produit environ 680 jetons de réponse. Pour un jeton écrit, il en a relu plus de
sept cents.

Le fournisseur ne publie pas l'énergie d'un jeton. On pondère donc les jetons par leur prix
catalogue, en supposant que le prix suit à peu près le calcul. C'est grossier : le prix
contient une marge et une stratégie commerciale. Avec cette pondération, sur la même
fenêtre, les lectures en cache font 70\,\% du total, les écritures en cache 20\,\% et les
réponses 10\,\%. Le coût d'une session tient donc moins à ce que le modèle écrit qu'à la
longueur de ce qu'il relit.

\section{Les chiffres du projet}

L'outil de l'assistant enregistre une ligne par appel, avec ses jetons, mais il efface ses
historiques après quatorze jours par défaut. Les enregistrements conservés couvraient, le
7 novembre 2026, la période du 24 septembre au 7 novembre : 44 des 92 jours du projet. Le
reste est extrapolé à partir de l'activité du dépôt (commits, lignes de transcription).
Pour passer des jetons au carbone, on encadre deux facteurs : l'énergie par jeton (de 0,5
à 5~mWh, d'après les chiffres publiés pour des modèles comparables) et l'intensité carbone
de l'électricité (de 0,2 à 0,4~kg\,CO$_2$e par kWh, faute de savoir où sont les centres
de données).

\begin{itemize}
  \item \textbf{L'inférence} domine tout. L'estimation centrale pour les trois mois est
        d'environ 110~kg\,CO$_2$e, dans une fourchette de 15~kg à 1,1~t. Ce sont des bornes
        obtenues en combinant les extrêmes de chaque facteur, et non des intervalles de
        confiance.
  \item \textbf{La fabrication du matériel} n'est pas dans ce chiffre. Un recoupement avec
        l'analyse du cycle de vie publiée par Mistral AI pour son modèle Mistral Large~2
        (\url{https://mistral.ai/news/our-contribution-to-a-global-environmental-standard-for-ai}),
        qui compte la fabrication des serveurs, donne cinq à neuf fois plus.
  \item \textbf{L'intégration continue} pèse peu. Au 7 novembre 2026, l'API de GitHub
        donnait la durée de 2\,769 tâches, pour 1\,541 exécutions : 271,5~heures en tout,
        dont 96\,\% pour le déploiement du site. Avec les coefficients publiés de la
        méthode Cloud Carbon Footprint, cela fait de 1,0 à 4,2~kg\,CO$_2$e.
  \item \textbf{L'eau} n'est pas estimée. Les deux chiffres publiés par réponse diffèrent de
        plus d'un facteur cent (0,26~mL pour la requête médiane de Gemini chez Google,
        45~mL pour une réponse de 400 jetons chez Mistral), et rien ne permet de choisir
        entre eux.
  \item \textbf{L'entraînement} des modèles est laissé de côté.
\end{itemize}

Pour donner une échelle, l'outil GES~1point5 trouve en moyenne 3,6~t\,CO$_2$e par an et
par membre de laboratoire en France (Mariette et al., 2022,
\url{https://doi.org/10.1088/2634-4505/ac84a4}). Les trois mois du projet en font quelques
pour cent au centre de la fourchette, et beaucoup plus si l'on compte le matériel. Ce
n'est pas rien. Ce n'est pas non plus ce qui décide, à lui seul, si un tel projet vaut
d'être fait.

\begin{regle}
Un chiffre de consommation porte sa date, la fenêtre qu'il couvre et la façon dont il a
été mesuré. Un chiffre extrapolé le dit, avec le facteur et ce qui le justifie.
\end{regle}

\section{Faire avec moins}

Le compte montre où agir. Presque tout se joue dans la longueur des contextes, un peu dans
le choix du modèle, très peu dans l'intégration continue.

\paragraph{Des sessions courtes.} Une session qui a lu dix fichiers relit les dix à chaque
tour, même quand le travail ne porte plus que sur l'un d'eux. On ouvre une session par
lot, on la ferme quand le lot est commité, et l'on donne à l'agent l'extrait dont il a
besoin plutôt que le fichier entier. Un sous-agent qui lit pour son compte et ne rend
qu'une conclusion garde le contexte principal court ; plusieurs sous-agents lancés pour
rien le rallongent d'autant.

\paragraph{Le plus petit modèle qui suffit.} Renommer des fichiers, relancer une
compilation, résumer une sortie de commande ne demandent pas le plus gros modèle. Le projet
l'a peu fait : sur la fenêtre mesurée, le modèle principal a reçu 15\,556 appels, le modèle
moyen 1\,584 et le petit 23, tous entre le 2 et le 7 novembre 2026. La marge est là.

\paragraph{Ne pas refaire ce qui n'a pas changé.} Le déploiement recompilait à chaque envoi
toutes les transcriptions en PDF, environ vingt-cinq minutes quand rien n'était en cache.
L'issue \#123 a mis chaque PDF en cache sous l'empreinte de ses sources, sauté le
déploiement pour les envois qui ne touchent pas au site et annulé les constructions
dépassées. Le 7 novembre 2026, les 32 derniers déploiements réussis avaient pris
2,5~minutes en moyenne. Le gain compte peu pour le carbone, puisque l'intégration continue
pèse quelques kilogrammes. Il compte pour le principe, qui est le même que pour le cache du
modèle : ne pas recalculer ce qu'on a déjà.

\paragraph{Un plafond par semaine.} L'abonnement de l'assistant donne une allocation
hebdomadaire. Le projet s'en fixe une part, suivie dans l'outil, et la dépasser demande
une décision. Un plafond ne mesure pas l'énergie, mais il oblige
à choisir ce qui passe en premier.

\paragraph{Des images à la résolution qui suffit.} Une image coûte des jetons à proportion
de sa taille. On lit une dactylographie à basse résolution et l'on garde les
agrandissements pour les formules et les passages douteux.

\paragraph{Un modèle plus petit, ou chez soi.} Un modèle ouvert hébergé par une université,
ou un modèle de taille moyenne sur une infrastructure qu'on connaît, permettrait de savoir
où l'électricité est produite et ce qu'elle émet. Le projet ne l'a pas essayé ; c'est la
suite logique de la règle qui garde tout le travail dans le dépôt, en texte brut, sans
dépendre d'un fournisseur.

\section{Ne pas demander au modèle}

La façon la plus sobre d'utiliser un modèle est de ne pas lui demander ce qu'un programme
fait mieux. Compter des pages, des théorèmes ou des lots se fait par un script, qui coûte
une fraction de seconde et donne le même résultat à chaque fois
(chapitre~\ref{ch:publier}). Une date se lit sur la page ou dans le catalogue ; on ne la
fait pas deviner. Un énoncé se vérifie dans Lean, qui répond oui ou non, plutôt qu'en
reposant la question au modèle jusqu'à ce qu'il réponde la même chose deux fois. Chaque
question qu'on ne pose pas est un contexte de plusieurs centaines de milliers de jetons
qu'on ne relit pas.

\section{L'effet rebond}

Rendre une activité moins chère tend à la faire croître. Pour l'entraînement des modèles,
Morand, Ligozat et Névéol trouvent que l'énergie et les impacts ont augmenté de façon
exponentielle en dix ans, malgré un matériel plus efficace
(\url{https://hal.science/hal-04839926v6}). Le projet en donne un exemple à sa petite
échelle. La lecture par la machine a été assez bon marché pour qu'on lise tout : 9\,414
pages en trois mois, presque deux fois les 5\,129 pages que les éditions faites à la main
avaient établies. Le calcul n'a pas remplacé un travail qui se serait fait autrement ; il a
servi à en faire davantage. Les 110~kg ne sont donc pas une économie : ils s'ajoutent.

On en tire deux règles de conduite. Un coût par page faible ne justifie pas un usage, car
c'est le total qui compte, et le total croît avec ce qu'on décide de lire. Le choix de ce
qu'on fait lire à la machine se fait pour lui-même, avant de lancer les lots : une première
lecture de tout un fonds a la valeur d'un inventaire, la même lecture refaite à chaque
reconstruction n'en a guère.

\section{Mesurer soi-même}

Les mesures de ce chapitre se refont avec deux sources. Les journaux de l'assistant donnent
les jetons par appel ; on allonge leur durée de conservation dès le début du projet, ce
que celui-ci n'a fait que le 7 novembre 2026, d'où l'extrapolation. L'API de l'intégration
continue donne la durée de chaque tâche. Un script dans le dépôt fait le compte, et
l'article cite la date du relevé.

\begin{terminal}
# durée de chaque tâche d'une exécution de l'intégration continue
gh api repos/<compte>/<dépôt>/actions/runs/<id>/jobs \
  --jq '.jobs[] | [.name, .started_at, .completed_at] | @tsv'
\end{terminal}

\begin{outils}
Les journaux de l'assistant pour les jetons ; \cmd{gh api} pour les durées de
l'intégration continue ; la méthode Cloud Carbon Footprint
(\url{https://www.cloudcarbonfootprint.org/docs/methodology}) pour convertir des heures de
machine ; GES~1point5 pour comparer à l'empreinte d'un laboratoire ; \texttt{actions/cache}
et un filtre de chemins pour ne pas tout reconstruire.
\end{outils}
